\ssr{ЗАКЛЮЧЕНИЕ}

В ходе выполнения курсовой работы была достигнута поставленная цель -- разработано программное обеспечение для визуализации трёхмерного огня методом систем частиц с учётом внешнего освещения и проведено исследование влияния ключевых параметров системы на производительность. Все поставленные задачи успешно выполнены.

\begin{enumerate}
	\item Проведён сравнительный анализ существующих методов визуализации и моделей освещения. На его основе выбрана и обоснована эмиссионная модель освещения, которая сочетает учёт самосвечения частиц с влиянием внешнего источника света через функцию затухания. Для решения задачи видимости выбран алгоритм z-буфера с поддержкой аддитивного смешивания, что обеспечивает корректное отображение полупрозрачных частиц и высокую производительность при большом количестве примитивов.
	
	\item Разработана структура программного обеспечения с чётким разделением ответственности между компонентами. Спроектированы основные классы для представления частицы, системы частиц, камеры, источника света и сцены, что обеспечило модульность и расширяемость системы.
	
	\item Реализован программный комплекс на C++ с использованием фреймворка Qt. Программа предоставляет интерактивный интерфейс, позволяющий пользователю в реальном времени создавать, настраивать и управлять несколькими независимыми системами частиц, а также изменять параметры внешнего освещения и виртуальной камеры.
	
	\item Проведено исследование производительности системы. Установлена закономерная зависимость частоты кадров от общего количества частиц в сцене. Проанализировано влияние внешнего освещения и алгоритма z-буфера на производительность. Результаты показывают, что выбранные алгоритмы обеспечивают приемлемую частоту кадров при визуализации сотен тысяч частиц, демонстрируя баланс между реалистичностью и быстродействием.
\end{enumerate}

Таким образом, разработанное программное обеспечение подтверждает работоспособность выбранного подхода и решает задачу интерактивной визуализации трёхмерного огня с учётом внешнего освещения. Результаты работы могут быть использованы в дальнейшем для оптимизации алгоритмов, внедрения более сложных моделей освещения или адаптации системы для моделирования других схожих явлений (дыма, взрывов, тумана).

\clearpage